\BourbakiBackSection{参考文献}

{[}1{]} C. F. GAUSS, Werke, Vol. 2, Göttingen, 1863.

{[}2{]} B. RIEMANN, Gesammelte mathematische Werke, 2nd edition, Leipzig (Teubner), 1892.

{[}3{]} B. RIEMANN, in Lettere di E. Betti a P. Tardy, Rend. Accad. Lincei (5) 24 \(^{1}\) (1915) pp.~517-519.

{[}4{]} G. CANTOR, Gesammelte Abhandlungen, Berlin (Springer), 1932.

{[}5{]} G. CANTOR and R. DEDEKIND, Briefwechsel, Actualités scientifiques et industrielles, no. 518, Paris (Hermann), 1937.

{[}6{]} A. SCHOENFLIES, Entwicklungen der Mengenlehre und ihrer Anwendungen, Part I, 2nd edition, Leipzig-Berlin (Teubner), 1913.

{[}7{]} E. BOREL, Leçons sur la théorie des fonctions, 2nd edition, Paris (Gauthier-Villars), 1914.

{[}8{]} G. Ascoli, Le curve limiti di una varietà data di curve, Mem. Accad. Lincei (3) 18 (1883) p.~521-586.

{[}9{]} J. HADAMARD, Sur certaines applications possibles de la théorie des ensembles, Verhandl. Intern. Math.-Kongress, Zürich, 1898, pp.~201-202.

{[}10{]} V. VOLTERRA, Theory of Functionals, London-Glasgow (Blackie and Son), 1930.

{[}11{]} D. HILBERT, Gesammelte Abhandlungen, vol.~3, Berlin (Springer), 1935, pp.~10-37 {[}= Jahresbericht des D.M.V., 8 (1900), p.~184, and Math. Ann., 59 (1904), p.~161{]}.

{[}12{]} D. HILBERT, Grundzüge einer allgemeinen Theorie der Integralgleichungen, 2nd edition, Leipzig-Berlin (Teubner), 1924.

{[}13{]} D. HILBERT, Grundlagen der Geometrie, 7th edition, Leipzig-Berlin (Teubner), 1930.

{[}14{]} E. SCHMIDT, Über die Auflösung linearer Gleichungen mit unendlich vielen Unbekannten, Rend. Palermo, 25 (1908), pp.~53-77.

{[}15{]} M. FRÉCHET, Sur quelques points du calcul fonctionnel, Rend. Palermo, 22 (1906), pp.~1-74.

{[}16{]} F. RIESZ, Stetigkeitsbegriff und abstrakte Mengenlehre, Atti del IV Congresso Intern. dei Matem., Bologna, 1908, vol.~2, pp.~18-24.

{[}17{]} F. HAUSDORFF, Grundzüge der Mengenlehre, 1st edition, Leipzig (Veit), 1914.

{[}18{]} E. H. MOORE and H. L. SMITH, A general theory of limits, Am. Journ. of Math., 44 (1922) pp.~102-121.

{[}19{]} A. TYCHONOFF, Über die topologische Erweiterung von Räumen, Math. Ann., 102 (1930), pp.~514-561.

{[}20{]} H. CARTAN, Théorie des filtres; filtres et ultrafiltres, C. R. Acad. Sc. Paris, 205 (1937), pp.~595-598 and pp.~777-779.

